%---------------------------------------------------------------------------%
%->> Backmatter
%---------------------------------------------------------------------------%
\chapter[致谢]{致\quad 谢}\chaptermark{致\quad 谢}% syntax: \chapter[目录]{标题}\chaptermark{页眉}
%\thispagestyle{noheaderstyle}% 如果需要移除当前页的页眉
%\pagestyle{noheaderstyle}% 如果需要移除整章的页眉

每个读到这段话的人，都是值得我感谢的人。\\

9岁的我梦想“成为一名科学家”，那时的我并不知道这意味着什么。在29岁时的回忆是不可靠的，因此我很难说清，童年时的自己究竟是对世界运行的规律怀着天然的好奇，还是仅仅觉得知道很多问题的答案会显得很厉害。不过可以确定的是，我并未实现9岁时的梦想。我离“科学家”这个词依然遥远，但离另一些东西的距离更近了。

“为什么要学神经科学？”是我读博期间被频繁问到的问题，在人工智能迅猛发展的今天尤其如此。我可以美化自己的答案，说：“我们不知道人工智能是否拥有意识；也不知道人工智能与人类智能是否能被同一套理论所解释”。当然我也可以呈现悲观的事实：“年轻的时候总是容易上当”。然后细数这个领域的圈地自萌与一地鸡毛，劝退每一个想要学习神经科学的人。有趣的是，这两种说辞都只在我一念之间。或许这个问题本就不存在确切的答案，任何回答都不过是一种“解释”。

那么，是否所有问题都没有答案呢？好消息是，这是一个自相矛盾的命题。答案，或者说某种真理，或许是存在的。但另一方面，我们并不知道人类能在多大程度上接近这种真理。科学的历史是一部犯错的历史。科学也是一种解释，是人类对真理的近似而非真理本身。如果科学也不可信，那我的博士生涯在追寻什么？这成为了我怀疑人生意义的顶点。

从任何意义上，我都不曾远离这份怀疑。但怀疑可以与生活并存。一个从未怀疑的人，和一个假装无事发生的人，可以过着类似的人生。这又回到了刚才所说的“一念之间”。于是我忽然理解了加缪说的：“重要的不是治愈，而是带着病痛活下去”，也理解了罗曼·罗兰所说的：“世界上只有一种真正的英雄主义，那就是在认清生活的真相后，依然热爱生活”。我并未成为某种英雄，也不会用“痛苦”来总结我的生活。但我意识到，有些人在没有意义的荒原上拣选了某些事物，并将其称为信仰。当一个人从信仰中获得力量时，信仰的内容变得无关紧要，是“信仰”这个行为本身带来了改变。

科学就可以是这样一种信仰，甚至近乎是一种宗教了。7年前的我可能会问自己：“我是否喜欢科学？”而在经历了这一切怀疑之后，我更想问的是：“我是否信仰科学？”我仍在迟疑，但事实上，我已经“几乎要信仰科学”了。我觉得科学与人生是类似的：没有人是带着既定的使命降生的，正如科学并非仅仅为了留下几则定论。科学的核心不是结论，是怀疑，更是在怀疑中前行；人生的核心不是成为“什么”，是“成为”，更是在“成为”一切的过程中悲喜交加的真实。科学从未提前拿到过剧本，是科学家们无数次的实践，将一个又一个抽象的可能具象化。同样，我不知道我的未来落在何处，这需要我在未来人生的无数次的选择中，一次次证明自己是什么样的人。

或许科学也可以被这样理解：它不是一种客观真理，也不是一种纯粹主观的解释，而是主客观相互作用中持续生成的经验。人生也是如此：我们并非先持有一个确凿的自我，再与世界打交道。恰恰相反，我们在与世界的交互中，逐渐获得了自己的存在。存在即是意义，如果非要感谢什么，那就感谢我们都恰好存在。\\

我知道这种祝福并无实质的意义，但我依然选择用它来结束我的博士生涯：

“祝我们，都能笃定地信仰着什么。”

\rightline{2026年12月}
\chapter{作者简历及攻读学位期间发表的学术论文与其他相关学术成果}

\section*{作者简历：}
2015年8月——2019年6月，在中国科学技术大学物理学院获得学士学位。

2019年9月——2026年7月，在中国科学院脑科学与智能技术卓越创新中心攻读博士学位。


\cleardoublepage[plain]% 让文档总是结束于偶数页，可根据需要设定页眉页脚样式，如 [noheaderstyle]
%---------------------------------------------------------------------------%
